\section{Problem 50: derivatives of the totient distribution function (Round 8)}

\subsection{1) ROUND-8 OBJECTIVE}

\textbf{Path (C) (obstruction / strategy barrier).}
Round~7 produced a dense $G_\delta$ set
\[
\Omega\subset (0,1)\setminus\mathcal E,
\qquad
\mathcal E:=\Bigl\{\frac{\varphi(n)}{n}:n\ge 1\Bigr\},
\]
on which all four one-sided Dini derivatives are maximally oscillatory:
\[
\underline D_-G(x)=\underline D_+G(x)=0,
\qquad
\overline D_-G(x)=\overline D_+G(x)=+\infty
\qquad(x\in\Omega).
\]
This still left one gap in the obstruction picture: Round~7 only asserted that one-sided secant slopes become arbitrarily small and arbitrarily large. It did \emph{not} yet identify the full set of intermediate slopes that actually occur.

The new objective is therefore to strengthen Round~7 from mere oscillation to \emph{exact slope realization}. We prove that on the same residual set $\Omega$, every positive slope occurs \emph{exactly} on both sides at arbitrarily small scales. Equivalently, the finite one-sided secant spectrum at every $x\in\Omega$ is all of $[0,\infty)$.

This is a genuine strengthening of Round~7: any hypothetical point where $G'(x)$ exists and is a finite positive number must avoid not merely a residual set of large/small oscillation, but a residual set on which \emph{every positive local secant slope is realized infinitely often on both sides}.

\subsection{2) Round-7 FOUNDATION USED}

We use the following earlier-round results as black boxes.

\begin{itemize}
\item \textbf{Round~1.} The limiting distribution function
\[
G(u):=\lim_{X\to\infty}\frac1X\#\Bigl\{n\le X:\frac{\varphi(n)}n\le u\Bigr\}
\qquad(0\le u\le 1)
\]
exists, is continuous and strictly increasing on $[0,1]$, and the associated measure $\mu=dG$ is purely singular with respect to Lebesgue measure. In particular,
\[
G'(x)=0\qquad\text{for Lebesgue-a.e. }x\in(0,1).
\]

\item \textbf{Round~3.} The value set
\[
\mathcal E:=\Bigl\{\frac{\varphi(n)}{n}:n\ge 1\Bigr\}
\]
is countable and dense in $[0,1]$, and every $t\in\mathcal E$ has infinite left derivative:
\[
\lim_{h\to 0^+}\frac{G(t)-G(t-h)}{h}=+\infty.
\]

\item \textbf{Round~7 main theorem.} There exists a dense $G_\delta$ set
\[
\Omega\subset (0,1)\setminus\mathcal E
\]
such that for every $x\in\Omega$ one has
\[
\underline D_-G(x)=\underline D_+G(x)=0,
\qquad
\overline D_-G(x)=\overline D_+G(x)=+\infty.
\]
Here
\[
\underline D_-G(x):=\liminf_{h\to0^+}\frac{G(x)-G(x-h)}{h},
\qquad
\overline D_-G(x):=\limsup_{h\to0^+}\frac{G(x)-G(x-h)}{h},
\]
\[
\underline D_+G(x):=\liminf_{h\to0^+}\frac{G(x+h)-G(x)}{h},
\qquad
\overline D_+G(x):=\limsup_{h\to0^+}\frac{G(x+h)-G(x)}{h}.
\]
\end{itemize}

\subsection{3) NEW INSIGHT / TOOL (ROUND-8)}

The new observation is very simple but powerful:

\begin{quote}
For fixed $x$, the one-sided secant quotients are continuous as functions of the scale $h$.
\end{quote}

Once Round~7 gives ``values below $c$'' and ``values above $c$'' arbitrarily close to $h=0$, continuity in $h$ forces an \emph{exact} solution of the equation ``secant slope $=c$'' at an intermediate scale. Doing this for every $c>0$ upgrades one-sided Dini oscillation to a complete local secant-spectrum theorem.

Concretely, we will prove:

\begin{quote}
\emph{For every $x\in\Omega$, every $c>0$, and every sufficiently small radius $\delta>0$, there exist }$h_-,h_+\in(0,\delta)$ \emph{such that}
\[
\frac{G(x)-G(x-h_-)}{h_-}=c,
\qquad
\frac{G(x+h_+)-G(x)}{h_+}=c.
\]
\end{quote}

Thus every positive slope is attained exactly, on both sides, at arbitrarily small scales around every $x\in\Omega$.

\subsection{4) ATTACK PLAN (ROUND-8)}

The exact Round~7 gap is that ``liminf $=0$ and limsup $=\infty$'' does not by itself state what happens at intermediate values. We close that gap as follows.

\begin{enumerate}
\item Fix $x\in\Omega$ and $c>0$.
\item Use Round~7 to find arbitrarily small scales where the left quotient is below $c$ and also arbitrarily small scales where it is above $c$.
\item Since the map $h\mapsto [G(x)-G(x-h)]/h$ is continuous on $(0,x)$, the intermediate value theorem gives a smaller scale with quotient exactly $c$.
\item Repeat on the right.
\item Package the result into a local ``secant spectrum'' theorem and derive an interval corollary: every nonempty open interval contains a point $x\notin\mathcal E$ at which any prescribed positive slope is realized on both sides, hence also by the whole secant through those two points.
\end{enumerate}

This rules out a new class of strategies: one cannot hope to locate a point of finite positive derivative by first excluding a particular local slope $c>0$ from nearby one-sided secants.

\subsection{5) WORK (ROUND-8)}

\subsubsection*{5.1. One-sided secant quotients as functions of the scale}

For $x\in(0,1)$ and $0<h<\min\{x,1-x\}$ define
\[
Q_h^-(x):=\frac{G(x)-G(x-h)}{h},
\qquad
Q_h^+(x):=\frac{G(x+h)-G(x)}{h}.
\]
Since $G$ is continuous, for each fixed $x$ the functions
\[
h\longmapsto Q_h^-(x)\qquad(0<h<x),
\qquad
h\longmapsto Q_h^+(x)\qquad(0<h<1-x)
\]
are continuous.

For $x\in(0,1)$ define the finite one-sided secant spectra
\[
\Sigma_-(x):=\Bigl\{L\in[0,\infty):\exists h_n\to0^+\text{ with }Q_{h_n}^-(x)\to L\Bigr\},
\]
\[
\Sigma_+(x):=\Bigl\{L\in[0,\infty):\exists h_n\to0^+\text{ with }Q_{h_n}^+(x)\to L\Bigr\}.
\]

\subsubsection*{5.2. Exact realization of every positive slope on the residual set}

\begin{theorem}[Exact one-sided secant-spectrum on the Round~7 residual set]\label{thm:round8-main}
Let $\Omega\subset(0,1)\setminus\mathcal E$ be the dense $G_\delta$ set furnished by Round~7, so that for every $x\in\Omega$,
\[
\underline D_-G(x)=\underline D_+G(x)=0,
\qquad
\overline D_-G(x)=\overline D_+G(x)=+\infty.
\]
Then for every $x\in\Omega$, every $c>0$, and every
\[
0<\delta<\min\{x,1-x\},
\]
there exist $h_-,h_+\in(0,\delta)$ such that
\[
Q_{h_-}^-(x)=\frac{G(x)-G(x-h_-)}{h_-}=c,
\qquad
Q_{h_+}^+(x)=\frac{G(x+h_+)-G(x)}{h_+}=c.
\]
Equivalently, for every such $x$ and $\delta$,
\[
(0,\infty)\subseteq\{Q_h^-(x):0<h<\delta\}
\qquad\text{and}\qquad
(0,\infty)\subseteq\{Q_h^+(x):0<h<\delta\}.
\]
In particular, for every $x\in\Omega$,
\[
\Sigma_-(x)=\Sigma_+(x)=[0,\infty).
\]
\end{theorem}

\begin{proof}
Fix $x\in\Omega$, $c>0$, and $0<\delta<\min\{x,1-x\}$.

We first treat the left side. Since $\underline D_-G(x)=0$, there exists $a\in(0,\delta)$ such that
\[
Q_a^-(x)<c.
\]
Since $\overline D_-G(x)=+\infty$, there exists $b\in(0,\delta)$ such that
\[
Q_b^-(x)>c.
\]
Let
\[
u:=\min\{a,b\},\qquad v:=\max\{a,b\}.
\]
Then $0<u\le v<\delta$, and the continuous function $h\mapsto Q_h^-(x)$ on $[u,v]$ takes one value below $c$ and one value above $c$. By the intermediate value theorem, there exists $h_-\in[u,v]\subset(0,\delta)$ such that
\[
Q_{h_-}^-(x)=c.
\]

The right side is identical. Since $\underline D_+G(x)=0$, choose $a'\in(0,\delta)$ with
\[
Q_{a'}^+(x)<c,
\]
and since $\overline D_+G(x)=+\infty$, choose $b'\in(0,\delta)$ with
\[
Q_{b'}^+(x)>c.
\]
Applying the intermediate value theorem to the continuous map $h\mapsto Q_h^+(x)$ on $[\min\{a',b'\},\max\{a',b'\}]$ yields $h_+\in(0,\delta)$ such that
\[
Q_{h_+}^+(x)=c.
\]

This proves the exact slope realization statement.

To identify the finite secant spectra, first note that for any $c>0$ and each integer $m\ge1$, applying the argument above with
\[
\delta_m:=\min\Bigl\{\frac1m,\frac{x}{2},\frac{1-x}{2}\Bigr\}
\]
produces $h_{-,m},h_{+,m}\in(0,\delta_m)$ with
\[
Q_{h_{-,m}}^-(x)=c,
\qquad
Q_{h_{+,m}}^+(x)=c.
\]
Hence $h_{\pm,m}\to0^+$ and so $c\in\Sigma_-(x)\cap\Sigma_+(x)$.

Next, because $\underline D_-G(x)=\underline D_+G(x)=0$, there exist sequences $u_m\to0^+$ and $v_m\to0^+$ with
\[
Q_{u_m}^-(x)\to0,
\qquad
Q_{v_m}^+(x)\to0,
\]
so $0\in\Sigma_-(x)\cap\Sigma_+(x)$.

Finally, by definition every element of $\Sigma_\pm(x)$ is nonnegative, so
\[
\Sigma_-(x)=\Sigma_+(x)=[0,\infty).
\]
\end{proof}


\subsubsection*{5.3. Interval realization of any prescribed positive slope}

\begin{corollary}[Every interval contains exact secants of every positive slope]\label{cor:interval-exact-slope}
Let $I\subset(0,1)$ be a nonempty open interval and let $c>0$. Then there exist points
\[
a<x<b,
\qquad x\in I\setminus\mathcal E,
\qquad [a,b]\subset I,
\]
such that
\[
\frac{G(x)-G(a)}{x-a}=c,
\qquad
\frac{G(b)-G(x)}{b-x}=c.
\]
Consequently,
\[
\frac{G(b)-G(a)}{b-a}=c.
\]
Thus every nonempty open interval contains a secant of exact slope $c$.
\end{corollary}

\begin{proof}
Because $\Omega$ is dense in $(0,1)$ and $\Omega\subset(0,1)\setminus\mathcal E$, choose
\[
x\in\Omega\cap I.
\]
Set
\[
\delta:=\tfrac12\min\{x-\inf I,\sup I-x\}>0.
\]
Then $0<\delta<\min\{x,1-x\}$ and
\[
[x-\delta,x+\delta]\subset I.
\]
Applying Theorem~\ref{thm:round8-main} with this $x$, the given $c$, and this $\delta$, we obtain $h_-,h_+\in(0,\delta)$ such that
\[
\frac{G(x)-G(x-h_-)}{h_-}=c,
\qquad
\frac{G(x+h_+)-G(x)}{h_+}=c.
\]
Define
\[
a:=x-h_-,\qquad b:=x+h_+.
\]
Then $a<x<b$ and $[a,b]\subset I$. Moreover,
\[
\frac{G(x)-G(a)}{x-a}=c,
\qquad
\frac{G(b)-G(x)}{b-x}=c.
\]
Adding the two identities gives
\[
G(b)-G(a)=[G(b)-G(x)]+[G(x)-G(a)]
=c(b-x)+c(x-a)=c(b-a),
\]
which is exactly
\[
\frac{G(b)-G(a)}{b-a}=c.
\]
\end{proof}

\subsubsection*{5.4. A sharper formulation of the obstruction}

Theorem~\ref{thm:round8-main} says that on a dense $G_\delta$ set disjoint from $\mathcal E$, the finite one-sided secant spectrum is already maximal:
\[
\Sigma_-(x)=\Sigma_+(x)=[0,\infty).
\]
So any point $x$ where $G'(x)$ exists and is a finite positive number must avoid a residual set on which \emph{every} positive local slope occurs infinitely often on each side.

This rules out another class of proof strategies. In particular, no argument can hope to prove the existence of a positive derivative by first isolating a neighborhood (or a large topological set of points) where some fixed local slope $c>0$ is \emph{forbidden} for nearby one-sided secants: Round~8 shows that every such slope occurs everywhere on a residual set.

\subsection{6) ADVERSARIAL VERIFICATION}

\begin{itemize}
\item \textbf{Does continuity in the scale $h$ really hold?}
Yes. For fixed $x$, the maps
\[
h\mapsto \frac{G(x)-G(x-h)}{h},
\qquad
h\mapsto \frac{G(x+h)-G(x)}{h}
\]
are quotients of continuous functions by the nonzero continuous function $h$, hence continuous on their natural domains.

\item \textbf{Do we truly get values both below and above $c$ inside every small interval $(0,\delta)$?}
Yes. The identities
\[
\underline D_\pm G(x)=0,
\qquad
\overline D_\pm G(x)=+\infty
\]
mean exactly that for every $\delta>0$ there are scales smaller than $\delta$ where the quotient is below $c$, and also scales smaller than $\delta$ where it is above $c$.

\item \textbf{What about the slope $c=0$?}
Exact equality $Q_h^\pm(x)=0$ is impossible for $h>0$ because $G$ is strictly increasing. Our theorem therefore states exact realization only for $c>0$. The value $0$ enters the spectrum as a \emph{limit point}, which is why
\[
\Sigma_\pm(x)=[0,\infty)
\]
rather than $(0,\infty)$.

\item \textbf{Why does the whole secant $[a,b]$ have slope $c$ in Corollary~\ref{cor:interval-exact-slope}?}
Because the left and right increments from the midpoint $x$ each have slope $c$, and these increments add:
\[
G(b)-G(a)=[G(b)-G(x)]+[G(x)-G(a)]=c(b-x)+c(x-a)=c(b-a).
\]

\item \textbf{Does this really go beyond Round~7?}
Yes. Round~7 gave only lower/upper Dini information. Round~8 identifies the full finite one-sided secant spectrum on a residual set and proves exact realization of every positive slope on both sides at arbitrarily small scales.
\end{itemize}

\subsection{7) FINAL (EXACTLY ONE)}

\textbf{UNRESOLVED (BUT STRICTLY ADVANCED).}

Round~8 strengthens Round~7 from one-sided oscillation to a full local secant-spectrum theorem. There exists a dense $G_\delta$ set
\[
\Omega\subset(0,1)\setminus\mathcal E
\]
such that for every $x\in\Omega$ and every $c>0$, the slope $c$ is realized exactly by arbitrarily small left secants and also by arbitrarily small right secants based at $x$. Equivalently,
\[
\Sigma_-(x)=\Sigma_+(x)=[0,\infty)
\qquad(x\in\Omega).
\]

As a consequence, every nonempty open interval contains secants of \emph{every} positive slope, and even contains points $x\notin\mathcal E$ from which a prescribed slope $c$ is realized on both sides simultaneously.

This is a strict advance beyond Round~7 and rules out another family of proof strategies. Any hypothetical point where $G'(x)$ exists and is a finite positive number must now avoid a residual set on which all positive one-sided local slopes occur infinitely often.

\subsection{8) COMPLETION ESTIMATE (MANDATORY)}

\textbf{COMPLETION: 91\%}

\subsection{9) REFERENCES}

\begin{enumerate}
\item I. J. Schoenberg, \emph{On asymptotic distributions of arithmetical functions}, Trans. Amer. Math. Soc. \textbf{39} (1936), 315--330.
\item P. Erd\H{o}s, \emph{On the distribution function of additive functions}, Ann. of Math. (2) \textbf{47} (1946), 1--20.
\item G. Tenenbaum and V. Toulmonde, \emph{Sur le comportement local de la r\'epartition de l'indicatrice d'Euler}, Funct. Approx. Comment. Math. \textbf{35} (2006), 321--338.
\item V. Toulmonde, \emph{Module de continuit\'e de la fonction de r\'epartition de $\varphi(n)/n$}, Acta Arith. \textbf{121} (2006), no.~4, 367--402.
\end{enumerate}
